%! Introducing Statistics: Do the Data We Collected Tell the Truth? — Unit 3, Lesson 3.1 — Student Packet Cover Sheet
\documentclass[10pt]{article}
\usepackage{apstats-article}
\usepackage{apstats-boxes}
\usepackage{ltablex}
\keepXColumns

\begin{document}

% Full-bleed navy banner
\begin{tikzpicture}[remember picture, overlay]
  \fill[navy] (current page.north west)
    rectangle ([yshift=-0.9in]current page.north east);
\end{tikzpicture}
\vspace{-0.8in}

\begin{center}
  {\color{white}\textbf{\LARGE AP Statistics}}\\[4pt]
  {\color{skymid}\large\textbf{Unit 3: Collecting Data}}\\[6pt]
  {\color{white}\textbf{\large Lesson 3.1 \quad Introducing Statistics: Do the Data We Collected Tell the Truth?}}
\end{center}

\vspace{0.22in}
\namedateperiod
\vspace{0.18in}

\begin{learningtargetbox}
\small
\begin{enumerate}[leftmargin=1.5em, itemsep=5pt]
  \item \ldots{} recognize that data collection methods not based on chance can produce
        untrustworthy conclusions.
  \item \ldots{} explain why a large non-random sample is not more trustworthy than a
        small random sample.
  \item \ldots{} identify whose opinions are represented --- and whose are missing ---
        in a given data collection scenario.
\end{enumerate}
\end{learningtargetbox}

\vspace{0.14in}

\begin{tocbox}
\small
\renewcommand{\arraystretch}{1.5}
\begin{tabularx}{\linewidth}{@{} c l X r @{}}
\rowcolor{sky}
\textbf{\#} & \textbf{Component} & \textbf{Description} & \textbf{Score} \\
\midrule
1 & Warm-Up       & Spiral review + chance vs.\ non-chance thinking   & \blank{1.2cm} \\
2 & Guided Notes  & Data collection methods; trustworthiness of data & \blank{1.2cm} \\
3 & Group Activity & Evaluating real-world data collection scenarios   & \blank{1.2cm} \\
4 & Exit Ticket   & Principal lunch-menu scenario                      & \blank{1.2cm} \\
5 & Homework      & Four non-random survey evaluations                 & \blank{1.2cm} \\
\midrule
  & \multicolumn{2}{r}{\textbf{Total}} & \blank{1.2cm} \\
\end{tabularx}
\end{tocbox}

\vspace{0.14in}

\begin{remindbox}
\small
The \textbf{most important idea in Unit 3}: the method used to collect data determines
what conclusions are valid.  Random chance is the foundation of trustworthy statistics.

\vspace{6pt}
\begin{tabularx}{\linewidth}{@{} >{\centering\arraybackslash\columncolor{goldbg}}p{0.06\linewidth}
                                    X @{}}
\textbf{\large!} &
\textbf{Big sample $\neq$ trustworthy sample.}\enspace
The 1936 Literary Digest poll contacted 10 million people and still got the
presidential election wrong --- because the sample was not randomly selected. \\[6pt]
\textbf{\large!} &
\textbf{Chance-based methods make conclusions valid.}\enspace
When every individual has an equal (or known) chance of being selected, the sample
is likely to represent the population. \\[6pt]
\textbf{\large!} &
\textbf{``Convenient'' data often tells us about convenient people.}\enspace
Voluntary respondents, friends, and online clickers are not random samples of the
population. Their answers reflect their opinions, not everyone's. \\
\end{tabularx}
\end{remindbox}

\end{document}
